\clearpage
\section{原文勘误与存疑清单（译者所加）}
\label{sec:errata}

下面列出翻译过程中发现的原文笔误、记号不一致与存疑之处。译文\textbf{一律照原文排印}，不改动数学内容；涉及理解的地方已在正文以译注说明。条目按原文出现顺序排列。带「\checkmark{}」者已对照原文源码与 PDF 双重确认。

\subsection*{一、图表与编号}
\begin{enumerate}[leftmargin=1.6em,itemsep=2pt]
\item （\checkmark{} §1b，原文 p.6）真下降方图中 $\Coh(Y)_{\infty}$ 向右的箭头标为 $i_{\infty}$、向下的箭头标为 $i_0$；这与 §1a（原文 p.5）的定义（$i_0$ 为零截面、$i_{\infty}$ 为无穷远截面）以及 §2b（原文 p.10）的真前推方图相反。按 §2b 与「$i_0$ 镜像于 $\cup_0$」的对应关系（原文 p.9--10），§1b 方图应读作 $i_0\colon\Coh(Y)\to\Coh(P(Y,\xi_2))$、$i_{\infty}\colon\Coh(Y)\to\Coh(P(Y,\xi_1))$。读者以正文的粘合约定 $\{z_1=\infty\}\cong\{z_2=0\}$ 为准（详见译注 9）。
\end{enumerate}

\subsection*{二、拼写与作者名}
\begin{enumerate}[leftmargin=1.6em,itemsep=2pt,resume]
\item （\checkmark{}）原文两处把作者名拼作 ``Chiang-Klem-Yau-Zaslow''（§1 开头与 p.5），应为 Chiang--Klemm--Yau--Zaslow（文献 \cite{CKYZ99}）。
\item （\checkmark{}§1a）``gi ve''（p.5）、``compatification''（p.4）、``firations''（p.5）、``sup posedly''（§1c）——均为排版时断词造成的拼写错误，译文按原意处理。
\item （\checkmark{}§6b）``gorup''（p.43）、``customly''（p.43）——应为 group、customarily。
\end{enumerate}

\subsection*{三、记号与公式（存疑，未改）}
\begin{enumerate}[leftmargin=1.6em,itemsep=2pt,resume]
\item （§5b，特征的定义）原文说特征取值于 $X$ 上平凡 $\mathbb{G}_m$-丛 $\mathcal{O}_X^{\times}$ 的整体截面，但紧接着用的是 $\Gamma(G,\mathcal{O}_G^{\times})$、$\mathcal{O}_{G\times G}^{\times}$；从上下文看应为 $G$ 上的 $\mathcal{O}_G^{\times}$。译文按 $G$ 处理并加译注。
\item （§5c，对角化群概形的定义）``the group $\cha$ is canonically identified with the character group $\widehat{G}$'' 中的 $\cha$ 应为 $\Lambda$。
\item （§5c，引理证明）出现 $\mathcal{O}_{G\times X}(G\times U)$-模的写法，与同段其余处的 $G\times X$ 不一致，其中 $U$ 疑为 $X$。
\item （§6b）某处下标写作 $i\in\mathbb{Z}$，据上下文应为 $i_0\in I$；证明末尾有一处多余的右括号。
\item （§6c）末式 $H^0\hom^{\bullet}\simeq\bigoplus_p\hom(F^p,G^p)$ 疑缺少双指标（应为双复形滤过分次的同调）。
\item （§5c）推论环境使用了标签 \texttt{thm:EmbeddingOfTheOrbitCategory}（标签名含 thm，环境其实是 cor）。仅为标签命名，不影响排版与编号，译文保持原标签以免破坏交叉引用。
\item （§6b）原文一处 ``\ldots{} invariant, Then we can define''，逗号后直接接大写 Then，应断句为句号；译文按句号处理。
\item （§6c）定理 6.10（扭 Künneth 公式）中 $\Ku^n$ 的定义漏写 $G,H$ 下标，按上下文应为 $\Ku^n_G(X)\otimes\Ku^n_H(Y)$ 一类带群下标的写法；译文照原文排印。
\item （§6c 推论证明）原文出现 $\Perf_{G\times G}(Y)$，按上下文应为 $\Perf_{G\times G}(X\times Y)$。
\item （§6b）原文以 quasi-coherent 描述 $\Coh_G(X;\mathfrak{U})_{\dg}$，与记号 $\Coh$（凝聚层）字面不符（系数在 Čech 复形中取拟凝聚层，整体仍是凝聚层导出范畴意义上的对象）；另 $\Res^G_H$ 与 $\Res_H^G$ 的上下标在原文中混用。译文统一保留各出现处的写法。
\item （§3a）原文 $\phi^d$ 的值域一处写作 $\hom_A$，按上下文应为 $\hom_B$。
\item （§3a）dg 范畴定义中 $\mu^k=0$ 的条件原文写作 $k\ge 2$，按 dg 范畴的通常定义应为 $k\ge 3$（$\mu^2$ 是一般不消失的合成）。
\item （§3b）某证明中 $p_{j+l+1}$ 未在上下文中定义，疑为 $p_{j+l}$ 或另一指标的笔误。
\end{enumerate}

\subsection*{四、附录 A 的笔误}
\begin{enumerate}[leftmargin=1.6em,itemsep=2pt]
\item （A.1 定义后的 display C）求和指标写作 $V_{j_0\cdots j_p}$，按上下标规则应为 $V_{j_0\cdots j_r}$ 一类；$\beta$ 的上标有一处误入下标（$f_{\beta\cdots^{\,k}}$）。
\item （display D）指标 $j_s$ 应为 $j_r$。
\item （链同伦证明）原文自指 \texttt{prp:STHd-relation}，该标签并不存在，应为 \texttt{prp:chain-homotopy-TS}；译文保留原引用并在此说明。
\end{enumerate}

\subsection*{五、译者的处理原则}
\begin{itemize}[leftmargin=1.6em,itemsep=2pt]
\item 原文笔误一律保留原样，只在译注或本表中指出；凡按原意「校正后翻译」之处，均在译注中写明所据。
\item 原文中被 \texttt{\%} 注释掉、未排入 PDF 的草稿图（如 §5a--§5c 的若干 tikzcd 图）不予补画，以免与发表版本不一致。
\item 定理、定义等环境的编号严格按原书计数（含 Remark、Example 等共享计数器），可与原文逐号对照；原文号见附录 \ref{sec:pagemap} 与各页注「原文 p.N」。
\end{itemize}
